% Einheit 3 von 4 – Gleichungen aufstellen und grafisch oder numerisch lösen (bank/gleichungen-loesen/e3.jsonl, zusammenbau v0.1)
%% TODO Zeitmarke Einheit 3: Marken-Zeile nicht lesbar
%% TODO Zweigzeile Teil 1: was hier gelernt wird (Satz in Schülersprache aus dem Einheitstitel) – Einheit 3
\einheitenkopf[e3]{Einheit 3 von 4 · Gleichungen aufstellen und grafisch oder numerisch lösen}
\zweigzeile{Abitur GK}
\setcounter{aufgabe}{16}

%% TODO Titel: Ich-kann-Satz fehlt in der Bank (Platzhalter: Kettenname) – Nr. 17
%% TODO Anweisung: ein Satz über der ersten Teilaufgabe fehlt in der Bank – Nr. 17
\begin{aufgabe}{Gleichungen aufstellen und grafisch oder numerisch lösen}
%% TODO \rechenplatz{n}: Schrittzahl des Grundfalls fehlt in der Bank (nur bei fester Schreibform, 2.3 b)
\begin{teile}
\teil Welche Gleichung? Die Höhe $h(t)$ eines Ballons erreicht 50 m. \kreuz{$h(t) = 50$} \\ \kreuz{$h(50) = t$} \\ \kreuz{$h'(t) = 50$}
\teil Der Graph zeigt die Höhe $h(t)$ einer Drohne ($t$ in s, $h$ in m). Wann ist sie 4 m hoch? Schreibe die Gleichung und lies ab.

\begin{ksys}[xmin=0,xmax=6,ymin=0,ymax=6,ablesen] \parabel{-1}{3}{5}{h} \end{ksys}
Gleichung: \leerfeld , t = \leerfeld
\teil Die Abbildung zeigt die Graphen von $h$ und $g$. Beschreibe, wie man die Lösungen von $h(x) - g(x) = 0$ am Bild findet, und lies alle ab. \hfill (Abitur 2019 GK) \\

\begin{ksys}[xmin=-4,xmax=4,ymin=-4,ymax=4,ablesen] \funktion{0.25*\x^3-2*\x}{h} \gerade{-1}{0}{g} \end{ksys}
x = \leerfeld
\teil Ein Tragseil wird durch $s$ beschrieben, 1 LE = 10 m. Zwei Punkte in der rechten Hälfte liegen 30 m waagerecht auseinander, der rechte 8 m höher. Gleichung für die $x$-Koordinate des höheren Punkts? \hfill (Abitur 2023 LK) \\ \leerfeld
\teil $a(x)$ ist die Zahl der Abonnenten eines Kanals in Hundert, $x$ in Wochen (Graph). Löse $a(x + 2) = a(x) + 6$ grafisch und deute die Lösung. \hfill (Abitur 2025 GK) \\

\begin{ksys}[xmin=0,xmax=6,ymin=0,ymax=10,ablesen] \parabel{0.5}{0}{0}{a} \end{ksys}
x = \leerfeld
\teil Die Wachstumsrate einer Pflanze ist $r(t) = 4t \cdot e^{-0{,}5t}$ ($t$ in Wochen, $r$ in cm pro Woche, Graph). Zu welchem Zeitpunkt ist die Rate genauso groß wie 2 Wochen davor? Löse mit dem Rechner (zwei Stellen) und markiere beide Punkte im Graphen. \hfill (Abitur 2026 LK) \\

\begin{ksys}[xmin=0,xmax=8,ymin=0,ymax=4] \funktion{4*\x*exp(-0.5*\x)}{r} \end{ksys}
t = \leerfeld
\end{teile}
\begin{gleichungsraster}[2]
\gl{\text{Die Leistung einer Solaranlage ist $L(x) = 3 \cdot \mathrm{sin}\left(\frac{\pi}{12} \cdot (x - 6)\right) + 1$ (kW; $x$ in Stunden seit 0 Uhr, $6 \le x \le 18$), das Maximum um 12 Uhr. Wann am Vormittag ist erstmals die Hälfte des Maximums erreicht? Gib die Uhrzeit in Stunden und Minuten an. (Abitur 2025 LK)}} &
\gl{\text{Der Wasserstand in einem Hafen ist $h(t) = 2 + \mathrm{sin}\left(\frac{\pi}{6} t\right)$ ($t$ in h, $h$ in m). Zu welchen Zeitpunkten mit $0 \le t \le 24$ beträgt er $2{,}5$ m? (Abitur 2022 LK)}} \\
\gl{\text{$f(x) = -x^2 + 6x - 4$. Bilde $f'(x)$ und löse $\frac{f(x) - 0}{x - 0} = f'(x)$ rechnerisch. (Abitur 2025 GK)}} &
\end{gleichungsraster}
\begin{teile}
\teil Die Abbildung zeigt den Graphen von $f(x) = 3 - 0{,}25 \cdot (x - 1)^2$. $P$ im I. und $Q$ im II. Quadranten liegen auf dem Graphen und sind jeweils die dem Ursprung $O$ diagonal gegenüberliegende Ecke eines achsenparallelen Quadrats. Skizziere beide Quadrate, gib Gleichungen für ihre Seitenlängen an und beschreibe, wie man das Verhältnis der Umfänge berechnet. \hfill (Abitur 2023 LK)

\begin{ksys}[xmin=-4,xmax=4,ymin=-1,ymax=4] \parabel{-0.25}{1}{3}{f} \end{ksys}
\end{teile}
\end{aufgabe}

%% TODO Titel: Ich-kann-Satz fehlt in der Bank (Platzhalter: Kettenname) – Nr. 18
%% TODO Anweisung: ein Satz über der ersten Teilaufgabe fehlt in der Bank – Nr. 18
\begin{aufgabe}{Schnittpunkte zweier Graphen mit dem Rechner ermitteln}
\begin{teile}
\teil Ein Teichrand wird durch $f(x) = -0{,}5x^2 + 4$ und $g(x) = e^{0{,}5x}$ modelliert (1 LE = 1 m). Bestimme mit dem Rechner die Koordinaten aller Schnittpunkte auf zwei Nachkommastellen. \hfill (Abitur 2018 LK) \\ S( \leerfeld | \leerfeld )
\end{teile}
\end{aufgabe}

%% TODO Titel: Ich-kann-Satz fehlt in der Bank (Platzhalter: Kettenname) – Nr. 19
%% TODO Anweisung: ein Satz über der ersten Teilaufgabe fehlt in der Bank – Nr. 19
\begin{aufgabe}{Gleichungen aufstellen und grafisch oder numerisch lösen – Fehler finden, Begründen, Anwendung, Darstellungswechsel}
\begin{teile}
\teil Bedingung: „Der Pegel $p(x)$ ist 4 Stunden später um 30 cm höher.“ Ein Schüler schreibt: \rechnung{p(4) &= p(0) + 30} Finde den Fehler und schreibe die richtige Gleichung.
\teil Der Rechner liefert für $\mathrm{sin}(x) = 0{,}3$ nur $x \approx 0{,}30$. Begründe, warum es für $0 \le x \le 2\pi$ eine weitere Lösung gibt.
\end{teile}
\begin{gleichungsraster}[2]
\gl{\text{Kerze A ist 20 cm hoch und brennt 2 cm pro Stunde ab, Kerze B ist 15 cm hoch und brennt 1 cm pro Stunde ab. Wann sind beide gleich hoch, und wie hoch?}} &
\end{gleichungsraster}
\begin{teile}
\teil Die Abbildung zeigt den Graphen von $f$ und die Gerade $y = 3$. Welche Gleichung wird an den Schnittstellen gelöst? Lies die Lösungen ab.

\begin{ksys}[xmin=-3,xmax=5,ymin=0,ymax=6,ablesen] \parabel{0.5}{1}{1}{f} \gerade{0}{3}{y=3} \end{ksys}
Gleichung: \leerfeld , x = \leerfeld
\end{teile}
\end{aufgabe}
